\section{AHSME 1970}
        \begin{enumerate}
        	\item (\textit{AHSME 1970}) The fourth power of $\sqrt{1+\sqrt{1+\sqrt{1}}}$ is:


 $\textbf{(A) }\sqrt{2}+\sqrt{3}\qquad \textbf{(B) }\frac{1}{2}(7+3\sqrt{5})\qquad \textbf{(C) }1+2\sqrt{3}\qquad \textbf{(D) }3\qquad \textbf{(E) }3+2\sqrt{2}$

 \textbf{Answer: }E



\textbf{Solution}

We can simplify the expression to $\sqrt{1+\sqrt{2}}$ Then we square it so it is now $1+\sqrt{2}$. We still have to square it one more time, getting $\boxed{(E)\ 3+2\sqrt{2}}$.

 


	\item (\textit{AHSME 1970}) A square and a circle have equal perimeters. The ratio of the area of the circle to the area of the square is:


 $\textbf{(A) }\frac{4}{\pi}\qquad \textbf{(B) }\frac{\pi}{\sqrt{2}}\qquad \textbf{(C) }\frac{4}{1}\qquad \textbf{(D) }\frac{\sqrt{2}}{\pi}\qquad \textbf{(E) }\frac{\pi}{4}$

 \textbf{Answer: }A



\textbf{Solution}

Let's say the circle has a circumference ( or perimeter of $4\pi$). Since the perimeter of the square is the same as the perimeter of the circle, the side length of the square is $\pi$. That means that the area of the square is $\pi^2$. The area of the circle is $4\pi$. So the area of the circle over the area of the square is  $\frac{4\pi}{\pi^2} \Rightarrow \frac{4}{\pi} \Rightarrow$ $\fbox{A}$



	\item (\textit{AHSME 1970}) If $x=1+2^p$ and $y=1+2^{-p}$, then $y$ in terms of $x$ is:


 $\textbf{(A) }\dfrac{x+1}{x-1}\qquad \textbf{(B) }\dfrac{x+2}{x-1}\qquad \textbf{(C) }\dfrac{x}{x-1}\qquad \textbf{(D) }2-x\qquad \textbf{(E) }\frac{x-1}{x}$

 \textbf{Answer: }C



\textbf{Solution}

Since we want the $y$ expression in terms of $x$, let's convert the $y$ expression. We can convert it to $1+ \frac{1}{2^p} \Rightarrow \frac{2^p+1}{2^p} \Rightarrow \frac{x}{x-1} \Rightarrow$ $\fbox{C}$



	\item (\textit{AHSME 1970}) Let $S$ be the set of all numbers which are the sum of the squares of three consecutive integers. Then we can say that:


 $\textbf{(A) }\text{No member of }S\text{ is divisible by }2\qquad\\ \textbf{(B) }\text{No member of }S\text{ is divisible by }3\text{ but some member is divisible by }11\qquad\\ \textbf{(C) }\text{No member of }S\text{ is divisible by }3\text{ or }5\qquad\\ \textbf{(D) }\text{No member of }S\text{ is divisible by }3\text{ or }7\qquad\\ \textbf{(E) }\text{None of these}$

 \textbf{Answer: }B



\textbf{Solution}

Consider $3$ consecutive integers $a, b,$ and $c$.  Exactly one of these integers must be divisible by 3; WLOG, suppose $a$ is divisible by 3.  Then $a \equiv 0 \pmod {3},  b \equiv 1 \pmod{3},$ and $c \equiv 2 \pmod{3}$.  Squaring, we have that $a^{2} \equiv 0 \pmod{3}, b^{2} \equiv 1 \pmod{3},$ and $c^{2} \equiv 1 \pmod{3}$, so $a^{2} + b^{2} + c^{2} \equiv 2 \pmod{3}$.  Therefore, no member of $S$ is divisible by 3.

 Now consider $3$ more consecutive integers $a, b,$ and $c$, which we will consider mod 11.  We will assign $k$ such that $a \equiv k \pmod{11}, b \equiv k + 1 \pmod{11},$ and $c \equiv k + 2 \pmod{11}$.  Some experimentation shows that when $k = 4, a \equiv 4 \pmod{11}$ so $a^{2} \equiv 5 \pmod{11}$.  Similarly, $b \equiv 5 \pmod{11}$ so $b^{2} \equiv 3 \pmod{11}$, and $c \equiv 6 \pmod{11}$ so $c^{2} \equiv 3 \pmod{11}$.  Therefore, $a^{2} + b^{2} + c^{2} \equiv 0 \pmod{11}$, so there is at least one member of $S$ which is divisible by 11.  Thus, $\fbox{B}$ is correct.

 


	\item (\textit{AHSME 1970}) If $f(x)=\dfrac{x^4+x^2}{x+1}$, then $f(i)$, where $i=\sqrt{-1}$, is equal to:


 $\textbf{(A) }1+i\qquad \textbf{(B) }1\qquad \textbf{(C) }-1\qquad \textbf{(D) }0\qquad  \textbf{(E) }-1-i$

 \textbf{Answer: }D



\textbf{Solution}

$i^4 = 1$ and $i^2=-1$, so the numerator is $0$. As long as the denominator is not $0$, which it isn't, the answer is $0 \Rightarrow$ $\fbox{D}$



	\item (\textit{AHSME 1970}) The smallest value of $x^2+8x$ for real values of $x$ is:


 $\textbf{(A) }-16.25\qquad \textbf{(B) }-16\qquad \textbf{(C) }-15\qquad \textbf{(D) }-8\qquad  \textbf{(E) }\text{None of these}$

 \textbf{Answer: }B



\textbf{Solution}

Let's imagine this as a quadratic equation. To find the minimum or maximum value, we always need to find the vertex of the quadratic equation. The vertex of the quadratic is $\frac{-b}{2a}$ in $ax^2+bx+c=0$. Then to find the output, or the y value of the quadratic, we plug the vertex "x" value back into the equation. In this quadratic, a=1, b=8, and c=0. So the "x" value is $\frac{-b}{2a} \Rightarrow \frac{-8}{2} = -4$. Plugging it back into $x^2 + 8x$, we get $16-32 = -16 \Rightarrow$ $\fbox{B}$



	\item (\textit{AHSME 1970}) Inside square $ABCD$ with side $s$, quarter-circle arcs with radii $s$ and centers at $A$ and $B$ are drawn. 
These arcs intersect at point $X$ inside the square. How far is $X$ from side $CD$?


 $\textbf{(A) }\frac{1}{2}s(\sqrt{3}+4)\qquad \textbf{(B) }\frac{1}{2}s\sqrt{3}\qquad \textbf{(C) }\frac{1}{2}s(1+\sqrt{3})\qquad\\ \textbf{(D) }\frac{1}{2}s(\sqrt{3}-1)\qquad  \textbf{(E) }\frac{1}{2}s(2-\sqrt{3})$

 \textbf{Answer: }E



\textbf{Solution}

All answers are proportional to $s$, so for ease, let $s=1$.

 Let $ABCD$ be oriented so that $A(0, 0), B(1, 0), C(1, 1), D(0, 1)$.

 The circle centered at $A$ with radius $1$ is $x^2 + y^2 = 1$.  The circle centered at $B$ with radius $1$ is $(x - 1)^2 + y^2 = 1$.  Solving each equation for $1 - y^2$ to find the intersection leads to $(x - 1)^2 = x^2$, which leads to $x = \frac{1}{2}$.

 Plugging that back in to $x^2 + y^2 = 1$ leads to $y^2 = 1 - \frac{1}{4}$, or $y = \pm \frac{\sqrt{3}}{2}$.  Since we want the intersection within the square where $0 < x, y < 1$, we take the positive solution, and the intersection is at $X(\frac{1}{2}, \frac{\sqrt{3}}{2})$.

 The distance from $X$ to side $CD$, which lies along the line $y=1$, is $1 - \frac{\sqrt{3}}{2}$, or $\frac{2  -\sqrt{3}}{2}$.  This is answer $\fbox{E}$ when $s=1$.

 


	\item (\textit{AHSME 1970}) If $a=\log_8{225}$ and $b=\log_2{15}$, then


 $\textbf{(A) }a=\frac{1}{2}b\qquad \textbf{(B) }a=\frac{2b}{3}\qquad \textbf{(C) }a=b\qquad \textbf{(D) }b=\frac{1}{2}a\qquad \textbf{(E) }a=\frac{3b}{2}$

 \textbf{Answer: }B



\textbf{Solution}

The solutions imply that finding the ratio $\frac{a}{b}$ will solve the problem.  We compute $\frac{a}{b}$, use change-of-base to a neutral base, rearrange the terms, and then use the reverse-change-of-base:

 $\frac{\log_8 225}{\log_2 15}$
$\frac{\frac{\ln 225}{\ln 8}}{\frac{\ln 15}{\ln 2}}$
$\frac{\ln 225 \ln 2}{\ln 15 \ln 8}$
$\frac{\ln 225}{\ln 15} \cdot \frac{\ln 2}{\ln 8}$
$\ln_{15} 225 \cdot \ln_8 2$

 Since $15^2 = 225$, the first logarithm is $2$.  Since $8^{\frac{1}{3}} = 2$, the second logarithm is $\frac{1}{3}$.

 Thus, we have $\frac{a}{b} = \frac{2}{3}$, or $a = \frac{2}{3}b$, which is option $\fbox{B}$.

 


	\item (\textit{AHSME 1970}) Points $P$ and $Q$ are on line segment $AB$, and both points are on the same side of the midpoint of $AB$.
Point $P$ divides $AB$ in the ratio $2:3$ and $Q$ divides $AB$ in the ratio $3:4$. If $PQ=2$, then the length of segment $AB$ is


 $\textbf{(A) }12\qquad \textbf{(B) }28\qquad \textbf{(C) }70\qquad \textbf{(D) }75\qquad  \textbf{(E) }105$

 \textbf{Answer: }C



\textbf{Solution}

In order, the points from left to right are $A, P, Q, B$.  Let the lengths between successive points be $x, 2, y$, respectively.  

 Since $\frac{AP}{PB} = \frac{2}{3}$, we have $\frac{x}{2 + y} = \frac{2}{3}$.

 Since $\frac{AQ}{QB} = \frac{3}{4}$, we have $\frac{x + 2}{y} = \frac{3}{4}$.

 The first equation gives $x = \frac{2}{3}(2 + y)$.  Plugging that in to the second equation gives:

 $\frac{\frac{2}{3}(2 + y) + 2}{y} = \frac{3}{4}$.

 $\frac{2}{3}(2 + y) + 2 = \frac{3}{4}y$
$\frac{4}{3} + \frac{2}{3}y + 2 = \frac{3}{4}y$
$\frac{10}{3} = \frac{1}{12}y$
$y = 40$

 Since $x = \frac{2}{3}(2 + y)$, we have $x = 28$.  The length of the entire segment is $x + 2 + y$, which is $70$, or option $\fbox{C}$.

 


	\item (\textit{AHSME 1970}) Let $F=.48181\cdots$ be an infinite repeating decimal with the digits $8$ and $1$ repeating. 
When $F$ is written as a fraction in lowest terms, the denominator exceeds the numerator by


 $\textbf{(A) }13\qquad \textbf{(B) }14\qquad \textbf{(C) }29\qquad \textbf{(D) }57\qquad  \textbf{(E) }126$

 \textbf{Answer: }D



\textbf{Solution}

Multiplying by $100$ gives $100F = 48.181818...$.  Subtracting the first equation from the second gives $99F = 47.7$, and all the other repeating parts cancel out.  This gives $F = \frac{47.7}{99} = \frac{477}{990} = \frac{159}{330} = \frac{53}{110}$.  Subtracting the numerator from the denominator gives $\fbox{D}$.

 


	\item (\textit{AHSME 1970}) If two factors of $2x^3-hx+k$ are $x+2$ and $x-1$, the value of $|2h-3k|$ is


 $\textbf{(A) }4\qquad \textbf{(B) }3\qquad \textbf{(C) }2\qquad \textbf{(D) }1\qquad  \textbf{(E) }0$

 \textbf{Answer: }E



\textbf{Solution}

From the Remainder Theorem, we have $2(-2)^3 - h(-2) + k = 0$ and $2(1)^3 - h(1) + k = 0$.  Simplifying both of those equations gives $-16 + 2h + k = 0$ and $2 - h + k = 0$.  Since $k = 16 - 2h$ and $k = h - 2$, we set those equal to get:

 $16 - 2h = h - 2$
$3h = 18$
$h = 6$

 This gives $k = 4$ when substituting back into either of the two equations, and $|2h - 3k| = 0$, which is answer $\fbox{E}$.

 


	\item (\textit{AHSME 1970}) A circle with radius $r$ is tangent to sides $AB$, $AD$, and $CD$ of rectangle $ABCD$ and passes through
the midpoint of diagonal $AC$. The area of the rectangle in terms of $r$, is


 $\textbf{(A) }4r^2\qquad \textbf{(B) }6r^2\qquad \textbf{(C) }8r^2\qquad \textbf{(D) }12r^2\qquad  \textbf{(E) }20r^2$

 \textbf{Answer: }C



\textbf{Solution}

$\fbox{C}$



	\item (\textit{AHSME 1970}) Given the binary operation $\ast$ defined by $a\ast b=a^b$ for all positive numbers $a$ and $b$. The for all positive $a,b,c,n$, we have


 $\textbf{(A) }a\ast b=b\ast a\qquad \textbf{(B) }a\ast (b\ast c)=(a\ast b)\ast c\qquad\\ \textbf{(C) }(a\ast b^n)=(a\ast n)\ast b\qquad \textbf{(D) }(a\ast b)^n=a\ast (bn)\qquad  \textbf{(E) }\text{None of these}$

 \textbf{Answer: }D



\textbf{Solution}

Let $a = 2, b = 3, c=n = 4$.  If all of them are false, the answer must be $E$.  If one does not fail, we will try to prove it.

 For option $A$, we have $2^3 = 3^2$, which is clearly false.

 For option $B$, we have $2^{81} = 8^{4}$, which is false.

 For option $C$, we have $2^{81} = 16^3$, which is false.

 For option $D$, we have $8^4 = 2^{12}$, which is true.

 The LHS is $(a^b)^n$.  By the elementary definition of exponentiation, this is $a^b$ multiplied by itself $n$ times.  Since each $a^b$ is actually $a$ multiplied $b$ times, the expression $(a^b)^n$ is $a$ multiplied by itself $bn$ times.  

 The RHS is $a^{bn}$.  This is $a$ multiplied by itself $bn$ times.

 Thus, the LHS is always equal to the RHS, so $\fbox{D}$ is the only correct statement.

 


	\item (\textit{AHSME 1970}) Consider $x^2+px+q=0$ where $p$ and $q$ are positive numbers. If the roots of this equation differ by $1$, then $p$ equals


 $\textbf{(A) }\sqrt{4q+1}\qquad \textbf{(B) }q-1\qquad \textbf{(C) }-\sqrt{4q+1}\qquad\\ \textbf{(D) }q+1\qquad \textbf{(E) }\sqrt{4q-1}$

 \textbf{Answer: }A



\textbf{Solution}

From the quadratic equation, the two roots of the equation are $\frac{-p\pm\sqrt{p^2-4q}}{2}$. The positive difference between these roots is $\sqrt{p^2 - 4q}$. Setting $\sqrt{p^2-4q}=1$ and isolating $p$ gives $\sqrt{4q+1}$, or choice $\boxed{\text{(A)}}$.

 


	\item (\textit{AHSME 1970}) Lines in the $xy$-plane are drawn through the point $(3,4)$ and the trisection points of the line segment 
joining the points $(-4,5)$ and $(5,-1)$. One of these lines has the equation 


 $\textbf{(A) }3x-2y-1=0\qquad \textbf{(B) }4x-5y+8=0\qquad \textbf{(C) }5x+2y-23=0\qquad\\ \textbf{(D) }x+7y-31=0\qquad  \textbf{(E) }x-4y+13=0$

 \textbf{Answer: }E



\textbf{Solution}

The trisection points of $(-4, 5)$ and $(5, -1)$ can be found by trisecting the x-coordinates and the y-coordinates separately.  The difference of the x-coordinates is $9$, so the trisection points happen at $-4 + \frac{9}{3}$ and $-4 + \frac{9}{3} + \frac{9}{3}$, which are $-1$ and $2$.  Similarly, the y-coordinates have a difference of $6$, so the trisections happen at $3$ and $1$.  So, the two points are $(-1, 3)$ and $(2, 1)$.

 We now check which line has both $(3, 4)$ and one of the two trisection points on it.  Plugging in $(x, y) = (3, 4)$ into all five of the equations works.  The point $(2, 1)$ doesn't work in any of the five lines.  However, $(-1, 3)$ works in line $E$.  Thus, the answer is $\fbox{E}$



	\item (\textit{AHSME 1970}) If $F(n)$ is a function such that $F(1)=F(2)=F(3)=1$, and such that $F(n+1)=\dfrac{F(n)\cdot F(n-1)+1}{F(n-2)}$ for $n\ge 3$,
then $F(6)$ is equal to


 $\textbf{(A) }2\qquad \textbf{(B) }3\qquad \textbf{(C) }7\qquad \textbf{(D) }11\qquad  \textbf{(E) }26$

 \textbf{Answer: }C



\textbf{Solution}

Plugging in $n=3$ gives $F(4) = \frac{F(3) \cdot F(2) + 1}{F(1)} = \frac{1 \cdot 1 + 1}{1} = 2$.

 Plugging in $n=4$ gives $F(5) = \frac{F(4) \cdot F(3) + 1}{F(2)} = \frac{2 \cdot 1 + 1}{1} = 3$.

 Plugging in $n=5$ gives $F(6) = \frac{F(5) \cdot F(4) + 1}{F(3)} = \frac{3 \cdot 2 + 1}{1} = 7$.

 Thus, the answer is $\fbox{C}$.

 


 \textbf{Sidenote} All the numbers in the sequence $F(n)$ are integers. In fact, the function $F$ satisfies $F(n)=4F(n-2)-F(n-4)$. (Prove it!).

 


	\item (\textit{AHSME 1970}) If $r\ge 0$, then for all $p$ and $q$ such that $pq\neq 0$ and $pr>qr$, we have


 $\textbf{(A) }-p>-q\qquad \textbf{(B) }-p>q\qquad \textbf{(C) }1>-q/p\qquad\\ \textbf{(D) }1<q/p\qquad \textbf{(E) }\text{None of These}$

 \textbf{Answer: }E



\textbf{Solution}

If $r>0$ and $pr>qr$, we can divide by the positive number $r$ and not change the inequality direction to get $p>q$.  Multiplying by $-1$ (and flipping the inequality sign because we're multiplying by a negative number) leads to $-p < -q$, which directly contradicts $A$.  Thus, $A$ is always false.

 If $p < 0$ (which is possible but not guaranteed), we can divide the true statement $p>q$ by $p$ to get $1 > \frac{q}{p}$.  This contradicts $D$.  Thus, $D$ is sometimes false, which is bad enough to be eliminated.

 If $(p, q, r) = (3, 2, 1)$, then the condition that $pr > qr$ is satisfied.  However, $-p = -3$ and $q = 2$, so $-p > q$ is false for at least this case, eliminating $B$.

 If $(p, q, r) = (2, -3, 1)$, then $pr > qr$ is also satisfied.  However, $-\frac{q}{p} = 1.5$, so $1 > -\frac{q}{p}$ is false, eliminating $C$.

 All four options do not follow from the premises, leading to $\fbox{E}$ as the correct answer.

 


	\item (\textit{AHSME 1970}) $\sqrt{3+2\sqrt{2}}-\sqrt{3-2\sqrt{2}}$ is equal to


 $\textbf{(A) }2\qquad \textbf{(B) }2\sqrt{3}\qquad \textbf{(C) }4\sqrt{2}\qquad \textbf{(D) }\sqrt{6}\qquad  \textbf{(E) }2\sqrt{2}$

 \textbf{Answer: }A



\textbf{Solution}

Square the expression:

 $(\sqrt{3+2\sqrt{2}}-\sqrt{3-2\sqrt{2}})^2=3+\sqrt{2}-2\sqrt{(3+\sqrt{2})(3-\sqrt{2})}+3-2\sqrt{2}=6-2\sqrt{9-8}=6-2\sqrt{1}=4$
$\Rightarrow\sqrt{3+2\sqrt{2}}-\sqrt{3-2\sqrt{2}}=\sqrt{4}=2\Rightarrow\boxed{A}$



	\item (\textit{AHSME 1970}) The sum of an infinite geometric series with common ratio $r$ such that $|r|<1$, is $15$,
and the sum of the squares of the terms of this series is $45$. The first term of the series is


 $\textbf{(A) }12\qquad \textbf{(B) }10\qquad \textbf{(C) }5\qquad \textbf{(D) }3\qquad  \textbf{(E) }2$

 \textbf{Answer: }C



\textbf{Solution}

We know that the formula for the sum of an infinite geometric series is $S = \frac{a}{1-r}$.

 So we can apply this to the conditions given by the problem. 

 We have two equations:

 
 \begin{align*} 15 &= \frac{a}{1-r} \\ 45 &= \frac{a^{2}}{1-r^{2}} \end{align*}


 We get

 
 \begin{align*} a &= 15 - 15r \\ a^{2} &= 45 - 45r^{2} \\ \\ (15 - 15r)^{2} &= 45 - 45r^{2} \\ 270r^{2} - 450r + 180 &= 0 \\ 3r^{2} - 5r + 2 &= 0 \\ (3r - 2)(r - 1) &= 0 \end{align*}


 Since $|r| < 1$, $r = \frac{2}{3}$, so plug this into the equation above and we get $a = 15 - 15r = 15 - 10 = \boxed{\textbf{(C)} \quad 5}$

 Solution by $\underline{\textbf{Invoker}}$



	\item (\textit{AHSME 1970}) Lines $HK$ and $BC$ lie in a plane. $M$ is the midpoint of line segment $BC$, and $BH$ and $CK$ are perpendicular to $HK$. Then we 


 $\textbf{(A) }\text{always have }MH=MK\qquad\\ \textbf{(B) }\text{always have }MH>BK\qquad\\ \textbf{(C) }\text{sometimes have }MH=MK\text{ but not always}\qquad\\ \textbf{(D) }\text{always have }MH>MB\qquad \\ \textbf{(E) }\text{always have }BH<BC$

 \textbf{Answer: }A



\textbf{Solution}

$\fbox{A}$



	\item (\textit{AHSME 1970}) On an auto trip, the distance read from the instrument panel was $450$ miles. 
With snow tires on for the return trip over the same route, the reading was $440$ miles. 
Find, to the nearest hundredth of an inch, the increase in radius of the wheels if the original radius was $15$ inches.


 $\textbf{(A) }.33\qquad \textbf{(B) }.34\qquad \textbf{(C) }.35\qquad \textbf{(D) }.38\qquad  \textbf{(E) }.66$

 \textbf{Answer: }B



\textbf{Solution}

$\fbox{B}$



	\item (\textit{AHSME 1970}) If the sum of the first $3n$ positive integers is $150$ more than the sum of the first $n$ positive integers,
then the sum of the first $4n$ positive integers is


 $\textbf{(A) }300\qquad \textbf{(B) }350\qquad \textbf{(C) }400\qquad \textbf{(D) }450\qquad  \textbf{(E) }600$

 \textbf{Answer: }A



\textbf{Solution 1}

We can setup our first equation as

 $\frac{3n(3n+1)}{2} = \frac{n(n+1)}{2} + 150$

 Simplifying we get

 $9n^2 + 3n = n^2 + n + 300 \Rightarrow 8n^2 + 2n - 300 = 0 \Rightarrow 4n^2 + n - 150 = 0$

 So our roots using the quadratic formula are

 $\dfrac{-b\pm\sqrt{b^2 - 4ac}}{2a} \Rightarrow \dfrac{-1\pm\sqrt{1^2 - 4\cdot(-150)\cdot4}}{2\cdot4} \Rightarrow \dfrac{-1\pm\sqrt{1+2400}}{8} \Rightarrow 6, -25/4$

 Since the question said positive integers, $n = 6$, so $4n = 24$
$\frac{24\cdot 25}{2} = 300$
$\fbox{A}$




\textbf{Solution 2}

Expressing as an equation:\begin{equation}\tag{1}\frac{3n(3n+1)}{2} = \frac{n(n+1)}{2} + 150.  \end{equation}

 The sum of the first 4n positive integers =\begin{equation}\tag{2}\frac{4n(4n+1)}{2}\end{equation}

 We will try to rearrange Equation (1) to give equation (2)

 $\frac{3n(3n+1)}{2} - \frac{n(n+1)}{2} = 150$
$=\frac{n(3(3n+1)-(n+1))}{2} = 150 =\frac{n(9n+3-n-1)}{2}$
$\frac{n(8n+2)}{2}= \frac{2n(4n+1)}{2}  = 150$
$\frac{4n(4n+1)}{2} = 2*150 = 300$

 300 is the answer

 $\fbox{A}$

 
$\sim$Melkor

 


	\item (\textit{AHSME 1970}) The number $10!$ ($10$ is written in base $10$), when written in the base $12$ system, ends in exactly $k$ zeroes. The value of $k$ is


 $\textbf{(A) }1\qquad \textbf{(B) }2\qquad \textbf{(C) }3\qquad \textbf{(D) }4\qquad  \textbf{(E) } 5$

 \textbf{Answer: }D



\textbf{Solution}

A number in base $b$ that ends in exactly $k$ zeros will be divisible by $b^k$, but not by $b^{k+1}$.  Thus, we want to find the highest $k$ for which $12^k | 10!$.

 There are $4$ factors of $3$:  $3, 6, 9$, and an extra factor from $9$.

 There are $8$ factors of $2$:  $2, 4, 6, 8, 10$, an extra factor from $4, 8$, and a third extra factor from $8$.

 So, $2^8 \cdot 3^4 = 4^4 \cdot 3^4 = 12^4$ will divide $10!$.  Thus, the answer is $\fbox{D}$.

 


	\item (\textit{AHSME 1970}) An equilateral triangle and a regular hexagon have equal perimeters. If the area of the triangle is $2$, then the area of the hexagon is


 $\textbf{(A) }2\qquad \textbf{(B) }3\qquad \textbf{(C) }4\qquad \textbf{(D) }6\qquad  \textbf{(E) }12$

 \textbf{Answer: }B



\textbf{Solution 1}

Let $ABCDEF$ be our regular hexagon, with centre $O$ - and join $AO, BO, CO, DO, EO,$ and $FO$. Note that we form six equilateral triangles with sidelength $\frac{s}{2}$, where $s$ is the sidelength of the triangle (since the perimeter of the two polygons are equal). If the area of the original equilateral triangle is $2$, then the area of each of the six smaller triangles is $1/4\times 2 = \frac{1}{2}$ (by similarity area ratios). 

 Thus, the area of the hexagon is $6\times \frac{1}{2}=3$, hence our answer is $\fbox{B}$



	\item (\textit{AHSME 1970}) For every real number $x$, let $[x]$ be the greatest integer less than or equal to $x$. 
If the postal rate for first class mail is six cents for every ounce or portion thereof, 
then the cost in cents of first-class postage on a letter weighing $W$ ounces is always


 $\textbf{(A) }6W\qquad \textbf{(B) }6[W]\qquad \textbf{(C) }6([W]-1)\qquad \textbf{(D) }6([W]+1)\qquad  \textbf{(E) }-6[-W]$

 \textbf{Answer: }E



\textbf{Solution}

This question is trying to convert the floor function, which is more commonly notated as $\lfloor x \rfloor$, into the ceiling function, which is $\lceil x \rceil$.  The identity is $\lceil x \rceil = -\lfloor -x \rfloor$, which can be verified graphically, or proven using the definition of floor and ceiling functions.

 However, for this problem, some test values will eliminate answers.  If $W = 2.5$ ounces, the cost will be $18$ cents.  Plugging in $W = 2.5$ into the five options gives answers of $15, 12, 6, 18, 18$.  This leaves options $D$ and $E$ as viable.  If $W = 2$ ounces, the cost is $12$ cents.  Option $D$ remains $18$ cents, while option $E$ gives $12$ cents, the correct answer.  Thus, the answer is $\fbox{E}$.

 


	\item (\textit{AHSME 1970}) The number of distinct points in the $xy$-plane common to the graphs of $(x+y-5)(2x-3y+5)=0$ and $(x-y+1)(3x+2y-12)=0$ is


 $\textbf{(A) }0\qquad \textbf{(B) }1\qquad \textbf{(C) }2\qquad \textbf{(D) }3\qquad  \textbf{(E) }4$

 \textbf{Answer: }B



\textbf{Solution 1}

The graph $(x + y - 5)(2x - 3y + 5) = 0$ is the combined graphs of $x + y - 5=0$ and $2x - 3y + 5 = 0$. Likewise, the graph $(x -y + 1)(3x + 2y - 12) = 0$ is the combined graphs of $x-y+1=0$ and $3x+2y-12=0$. All these lines intersect at one point, $(2,3)$.Therefore, the answer is $\fbox{(B) 1}$.

 



\textbf{Solution 2}

We need to satisfy both $(x+y-5)(2x-3y+5)=0$ and $(x-y+1)(3x+2y-12)=0$. In order to do this, let us look at the first equation. Either $x+y=5,$ or $2x-3y=-5.$ For the second equation, either $x-y=-1,$ or $3x+2y=12.$ Thus, we need to solve 4 systems of equations: $x+y=5$ and $x-y=-1$, $x+y=5$ and $3x+2y=12$, $2x-3y=-5$ and $x-y=-1$, and finally $2x-3y=-5$ and $3x+2y=12.$ Solving all of these systems of equations is pretty trivial, and all of them come out to be $(2,3).$ Thus, they only intersect at $1$ point, and our answer is $\fbox{(B) 1}$.

 


	\item (\textit{AHSME 1970}) In a triangle, the area is numerically equal to the perimeter. What is the radius of the inscribed circle?


 $\textbf{(A) }2\qquad \textbf{(B) }3\qquad \textbf{(C) }4\qquad \textbf{(D) }5\qquad  \textbf{(E) }6$

 \textbf{Answer: }A



\textbf{Solution}

One of the most common formulas involving the inradius of a triangle is $A = rs$, where $A$ is the area of the triangle, $r$ is the inradius, and $s$ is the semiperimeter.

 The problem states that $A = p = 2s$.  This means $2s = rs$, or $r = 2$, which is option $\fbox{A}$.

 


	\item (\textit{AHSME 1970}) In triangle $ABC$, the median from vertex $A$ is perpendicular to the median from vertex $B$. 
If the lengths of sides $AC$ and $BC$ are $6$ and $7$ respectively, then the length of side $AB$ is


 $\textbf{(A) }\sqrt{17}\qquad \textbf{(B) }4\qquad \textbf{(C) }4\dfrac{1}{2}\qquad \textbf{(D) }2\sqrt{5}\qquad \textbf{(E) }4\frac{1}{4}$

 \textbf{Answer: }A



\textbf{Solution}

$\fbox{A}$let the midpoint be M,N ( i.e. AM,BN are the medians); connecting MN we know that AB = 2x and MN = x hence apply stewart's theorem in triangle ABC with median MN first and then apply stewart's in triangle BNC with median MN

 


	\item (\textit{AHSME 1970}) It is now between $10:00$ and $11:00$ o'clock, and six minutes from now, the minute hand of the 
watch will be exactly opposite the place where the hour hand was three minutes ago. What is the exact time now?


 $\textbf{(A) }10:05\dfrac{5}{11}\qquad \textbf{(B) }10:07\dfrac{1}{2}\qquad \textbf{(C) }10:10\qquad\\ \textbf{(D) }10:15\qquad \textbf{(E) }10:17\dfrac{1}{2}$

 \textbf{Answer: }D



\textbf{Solution}

$\fbox{D}$50 + (m-3)/12 = 30 + (m+6)= 10hours + m = 10: 15

 


	\item (\textit{AHSME 1970}) 
                    \begin{figure}[H]
                        \centering
                        \includegraphics[width=0.5\textwidth]{Problem/AHSME_I_1970_30_Problem_Figure_1.png}
                    \end{figure}
                    


 In the accompanying figure, segments $AB$ and $CD$ are parallel, 
the measure of $\angle{D}$ is twice the measure of $\angle{B}$, and the measures of segments $AD$ and $CD$
are $a$ and $b$ respectively. Then the measure of $AB$ is equal to


 $\textbf{(A) }\dfrac{1}{2}a+2b\qquad \textbf{(B) }\dfrac{3}{2}b+\dfrac{3}{4}a\qquad \textbf{(C) }2a-b\qquad \textbf{(D) }4b-\frac{1}{2}a\qquad \textbf{(E) }a+b$

 \textbf{Answer: }E



\textbf{Solution}

With reference to the diagram above, let $E$ be the point on $AB$ such that $DE||BC$. Let $\angle ABC=\alpha$. We then have $\alpha =\angle AED = \angle EDC$ since $AB||CD$, so $\angle ADE=\angle ADC-\angle BDC=2\alpha-\alpha = \alpha$, which means $\triangle AED$ is isosceles. 

 Therefore, $AB=AE+EB=a+b$, hence our answer is $\fbox{E}$.

 


	\item (\textit{AHSME 1970}) If a number is selected at random from the set of all five-digit numbers in which the sum of the digits is equal to $43$,
what is the probability that this number is divisible by $11$?


 $\textbf{(A) }2/5\qquad \textbf{(B) }1/5\qquad \textbf{(C) }1/6\qquad \textbf{(D) }1/11\qquad  \textbf{(E) }1/15$

 \textbf{Answer: }B



\textbf{Solution}

For the sums of the digits of be equal to 43, this means that the number's digits must fall under one of two categories. It either has 4 9's and a 7, or 3 9's and 2 8's. You can use casework from there to find that there are 15 5-digit numbers in which the sum of the digits is equal to 43. To find the number of ways that the number is divisible by 11, use the divisibility rule for 11, letting the number be represented by $abcde$. This gives 3 possible numbers that are divisible by 11, so the answer is $\frac{3}{15}=\frac{1}{5}=\fbox{B}$



	\item (\textit{AHSME 1970}) $A$ and $B$ travel around a circular track at uniform speeds in opposite directions, 
starting from diametrically opposite points. If they start at the same time, 
meet first after $B$ has travelled $100$ yards, and meet a second time $60$ yards before $A$ completes one lap, 
then the circumference of the track in yards is


 $\textbf{(A) }400\qquad \textbf{(B) }440\qquad \textbf{(C) }480\qquad \textbf{(D) }560\qquad  \textbf{(E) }880$

 \textbf{Answer: }C



\textbf{Solution}

Let $x$ be half the circumference of the track. They first meet after $B$ has run $100$ yards, meaning that in the time $B$ has run $100$ yards, $A$ has run $x-100$ yards. The second time they meet is when $A$ is 60 yards before he completes the lap. This means that in the time that $A$ has run $2x-60$ yards, $B$ has run $x+60$ yards. Because they run at uniform speeds, we can write the equation
 \[\frac{100}{x-100}=\frac{x+60}{2x-60} .\]Cross multiplying,
 \[200x-6000=x^2-40x-6000\]Adding $6000$ to both sides and simplifying, we have
 \[200x=x^2-40x\]
 \[240x=x^2\]
 \[x=240.\]Because $x$ is only half of the circumference of the track, the answer we are looking for is $2 \cdot 240 = 480, \text{or } \fbox{C}$.

 


	\item (\textit{AHSME 1970}) Find the sum of the digits of all numerals in the sequence $1,2,3,4,\cdots ,10000$.


 $\textbf{(A) }180,001\qquad \textbf{(B) }154,756\qquad \textbf{(C) }45,001\qquad \textbf{(D) }154,755\qquad  \textbf{(E) }270,001$

 \textbf{Answer: }A



\textbf{Solution 1}

We can find the sum using the following method. We break it down into cases. The first case is the numbers $1$ to $9$. The second case is the numbers $10$ to $99$. The third case is the numbers $100$ to $999$. The fourth case is the numbers $1,000$ to $9,999$. And lastly, the sum of the digits in $10,000$. The first case is just the sum of the numbers $1$ to $9$ which is, using $\frac{n(n+1)}{2}$, $45$. In the second case, every number $1$ to $9$ is used $19$ times. $10$ times in the tens place, and $9$ times in the ones place. So the sum is just $19(45)$. Similarly, in the third case, every number $1$ to $9$ is used  $100$ times in the hundreds place, $90$ times in the tens place, and $90$ times in the ones place, for a total sum of $280(45)$. By the same method, every number $1$ to $9$ is used $1,000$ times in the thousands place, $900$ times in the hundreds place, $900$ times in the tens place, and $900$ times in the ones place, for a total of $3700(45)$. Thus, our final sum is $45+19(45)+280(45)+3700(45)+1=4000(45)+1=\boxed{\text{A)}180001}.$




\textbf{Solution 2}

Consider the numbers from $0000-9999$. We have $40000$ digits and each has equal an probability of being $0,1,2....9$.Thus, our desired sum is $4000\left( \frac{9 \cdot 10}{2} \right)+1=4000(45)+1=\boxed{\text{A)}180001}.$

 \textit{Credit: Math1331Math}

 



\textbf{Solution 3}

As in Solution 2, we consider the four digit numbers from $0000-9999.$ We see that we have $10000\times4=40000$ digits with each digit  appearing equally.

 Thus, the digit sum will be the average of the digits multiplied by $40000.$ The digit average comes out to be $\frac{0+9}{2},$ since all digits are consecutive.

 So, our answer will be $\frac{9}{2} \times 40000 = 180000.$ However, since we purposely did not include $10000,$ we add one to get our final answer as $\boxed{\textbf{(A)}180001}.$

 Solution by davidaops.

 \textit{Credit to Math1331Math}

 


	\item (\textit{AHSME 1970}) The greatest integer that will divide $13511, 13903$, and $14589$ and leave the same remainder is


 $\textbf{(A) }28\qquad \textbf{(B) }49\qquad \textbf{(C) }98\qquad\\ \textbf{(D) }\text{an odd multiple of }7\text{ greater than }49\qquad\\ \textbf{(E) }\text{an even multiple of }7\text{ greater than }98$

 \textbf{Answer: }C



\textbf{Solution}

We know that 13903 minus 13511 is equivalent to 392. Additionally, 14589 minus 13903 is equivalent to 686. Since we are searching for the greatest integer that divides these three integers and leaves the same remainder, the answer resides in the greatest common factor of 686 and 392. Therefore, the answer is 98, or $\fbox{C}$

 "Credit: Skuppi3"

 


	\item (\textit{AHSME 1970}) A retiring employee receives an annual pension proportional to the square root of the number of years of his service. 
Had he served $a$ years more, his pension would have been $p$ dollars greater, whereas, had he served $b$ years more $b\neq a$, 
his pension would have been $q$ dollars greater than the original annual pension. Find his annual pension in terms of $a,b,p$, and $q$.


 $\textbf{(A) }\dfrac{p^2-q^2}{2(a-b)}\qquad \textbf{(B) }\dfrac{(p-q)^2}{2\sqrt{ab}}\qquad \textbf{(C) }\frac{ap^2-bq^2}{2(ap-bq)}\qquad \textbf{(D) }\frac{aq^2-bp^2}{2(bp-aq)}\qquad \textbf{(E) }\sqrt{(a-b)(p-q)}$

 \textbf{Answer: }D



\textbf{Solution}

Note the original pension as $k\sqrt{x}$, where $x$ is the number of years served. Then, based on the problem statement, two equations can be set up. 

 
 \[k\sqrt{x+a} = k\sqrt{x} + p\]
 \[k\sqrt{x+b} = k\sqrt{x} + q\]

 Square the first equation to get

 
 \[k^2x + ak^2 = k^2x + 2pk\sqrt{x} + p^2.\]Because both sides have 
 \[k^2x\], they cancel out. Similarly, the second equation will become $bk^2 = q^2 + 2qk\sqrt{x}.$ Then, $a$ can be multiplied to the second equation and $b$ can be multiplied to the first equation so that the left side of both equations becomes $abk^2$. Finally, by setting the equations equal to each other, 
 \[bp^2 + 2bpk\sqrt{x} = aq^2 + 2aqk\sqrt{x}.\] Isolate $k\sqrt{x}$ to get $\fbox{D} = \frac{aq^2-bp^2}{2(bp-aq)}$.

 


\end{enumerate}